% Copy to content/blog/<slug>.tex — the file name is the URL (/blog/<slug>).
% This one post shows every feature of the blog; delete what yours does not use.
% Put images in public/blog/<slug>/ and reference them from the site root.
% Everything before \begin{document} is metadata.
%
% The same file compiles with latexmk (mitaa.sty supplies the site's macros):
%   npm run blog:check <slug>              compare the site's numbering with LaTeX's
%   npm run blog:check <slug> -- --publish also copy the PDF to public/docs/blog/
\documentclass[11pt,a4paper]{article}
% \usepackage[dutch]{babel}  % Dutch labels on the site and in the PDF; English without it
\usepackage{amsmath, booktabs, graphicx}
\usepackage{mitaa}

\title{Post Title}
\subtitle{One line under the title, in italics}
\author{Mitchel Reints}
\affiliation{Independent, Utrecht}
\date{2026-01-01}
\version{v1}{}  % \version{v2}{2026-02-01} shows "v2 · revised …"
\kind{Report}  % the small label above the title
\keywords{embedded, control, measurement}
\excerpt{One or two sentences for the blog index card and the meta description.}
\tags{Embedded, C}
\coverimage{/blog-example.png}  % your own: /blog/<slug>/cover.png (index card and social image)
\xpost{https://x.com/handle/status/0000000000000000000}
% \pdf{/docs/blog/<slug>.pdf}  % a "Download as PDF" button; blog:check --publish puts it there
\paginate{auto}  % auto (6+ sections and 2000+ words) | always | never. \clearpage sets pages yourself, as below
% \draft  % renders under npm run dev only

\begin{document}
\maketitle

% An abstract before the first \section goes into the header, above the text,
% with the keywords under it.
\begin{abstract}
The question, the method and the answer, in four or five sentences.
\end{abstract}

\section{Text and structure}
\label{sec:text}

Open with the problem in two or three sentences: what you were trying to do, and why it was harder than it looked. Text can be \textbf{bold}, \emph{emphasised}, \textbf{\emph{both}} or \texttt{code}, and a \href{https://example.com}{link} takes any address. The characters LaTeX reserves are typed with a backslash: 50\,\% of \$5 \& 3\_4 \#1.

Each \texttt{\textbackslash{}section} is numbered and listed in the table of contents. Every \texttt{\textbackslash{}clearpage} in the source starts a new page on the site, and a page may hold several sections, as this one does: it is called after its first section, and the page links below the text step through them.

\subsection{Lists}

\begin{itemize}
  \item \texttt{itemize} for points without an order;
  \item lists can nest:
    \begin{itemize}
      \item one level down;
      \item and another item.
    \end{itemize}
  \item back at the top.
\end{itemize}

\begin{enumerate}
  \item State the problem.
  \item Describe what you tried.
  \item Give the result.
\end{enumerate}

\subsubsection{Quotes and footnotes}

A quotation sets its attribution flush right underneath.

\begin{quote}[--- Someone worth quoting]
A quotation, in a few lines.
\end{quote}

Footnotes are numbered through the whole post and printed at the foot of the page.\footnote{Measured at room temperature.}

\clearpage

\section{Callouts}
\label{sec:callouts}

Five kinds, each with an optional title in square brackets. The kind is shown by the colour of the rail.

\begin{note}
A \texttt{note} for background the reader may want.
\end{note}

\begin{tip}[A title]
A \texttt{tip} that saves time.
\end{tip}

\begin{important}
An \texttt{important} point to keep in mind.
\end{important}

\begin{warning}
A \texttt{warning} about something that goes wrong quietly.
\end{warning}

\begin{caution}
A \texttt{caution} before something that cannot be undone.
\end{caution}

\section{Code}
\label{sec:code}

\texttt{lstlisting} takes the language and a file name; \texttt{minted} the same, as \texttt{filename=}. A \texttt{verbatim} block is printed without colour.

\begin{lstlisting}[language=c, title=control.c]
void control_loop(void) {
    float error = setpoint - read_adc();
    set_pwm(kp * error);
}
\end{lstlisting}

Give a listing a \texttt{caption=} and a \texttt{label=} and it is numbered like a figure, so \Cref{lst:pid} can point at it.

\begin{lstlisting}[language=c, caption={The integral term with anti-windup}, label=lst:pid]
integral += error * dt;
integral = clamp(integral, -I_MAX, I_MAX);
\end{lstlisting}

\begin{minted}[filename=blink.py]{python}
led = Pin(25, Pin.OUT)
while True:
    led.toggle()
    time.sleep(0.5)
\end{minted}

\begin{verbatim}
$ make flash
avrdude: 1024 bytes of flash written
\end{verbatim}

\clearpage

\section{Mathematics}
\label{sec:math}

All of KaTeX is available. Inline maths goes between dollar signs, like $K_p$ and $x \in \R^n$, and display maths between \texttt{\textbackslash[} and \texttt{\textbackslash]}:
\[
  \int_0^\infty e^{-x^2} \dd x = \frac{\sqrt{\pi}}{2}
\]

A numbered \texttt{equation} takes the next number in the post and can be pointed at with \texttt{\textbackslash{}eqref}: the controller is \eqref{eq:pi}.
\begin{equation}
u(t) = K_p \, e(t) + K_i \int_0^t e(\tau)\, d\tau
\label{eq:pi}
\end{equation}

\subsection{Alignment, matrices and cases}

\[
\begin{aligned}
  \nabla \cdot \vect{E} &= \frac{\rho}{\eps_0} \\
  \nabla \cdot \vect{B} &= 0
\end{aligned}
\qquad
A = \begin{pmatrix} 1 & 2 \\ 3 & 4 \end{pmatrix}
\qquad
\abs{x} = \begin{cases} x & x \ge 0 \\ -x & x < 0 \end{cases}
\]

\subsection{Chemistry and units}

Reactions are written with \texttt{\textbackslash{}ce\{\}}: $\ce{H2O}$ and
\[
  \ce{2H2 + O2 -> 2H2O}
\]
The site also defines \texttt{\textbackslash{}R \textbackslash{}C \textbackslash{}Z \textbackslash{}N \textbackslash{}eps \textbackslash{}dd \textbackslash{}abs \textbackslash{}norm \textbackslash{}set \textbackslash{}vect \textbackslash{}unit}, as in $v = 3\unit{m/s}$ and $\norm{\vect{v}} < 1$; \texttt{mitaa.sty} gives them the same meaning in the PDF.

\section{Theorems and proofs}
\label{sec:theorems}

Theorems, lemmas, propositions, corollaries and definitions share one counter; examples have their own. A \texttt{proof} is not numbered.

\begin{definition}[Settling time]
The time after which the output stays within 2\,\% of its final value.
\end{definition}

\begin{theorem}[Stability]
\label{thm:stable}
The closed loop is stable for every $K_p > 0$ and $K_i \ge 0$.
\end{theorem}

\begin{proof}
Both poles of the characteristic polynomial lie in the left half-plane.
\end{proof}

\begin{lemma}
A lemma, to be used in the proof of something larger.
\end{lemma}

\begin{proposition}
A proposition.
\end{proposition}

\begin{corollary}
A corollary of \Cref{thm:stable}.
\end{corollary}

\begin{example}
\Cref{thm:stable} still says nothing about how fast the loop settles.
\end{example}

\clearpage

\section{Figures, tables and charts}
\label{sec:figures}

Figures, tables, equations and theorems are numbered as LaTeX numbers them, so a reference can sit anywhere, even before the thing it points at.

\subsection{Figures}

\begin{figure}
\includegraphics[width=0.6\textwidth]{/blog-example.png}  % your own: /blog/<slug>/diagram.png
\caption{What the reader should notice in this picture}
\label{fig:design}
\end{figure}

\subsection{Tables}

A caption written above the table is printed above it; below the content, below it. \texttt{booktabs} rules and plain \texttt{\textbackslash{}hline}s both work, and the first row becomes the header.

\begin{table}
\caption{Main parts}
\label{tab:parts}
\begin{tabular}{llr}
\toprule
Part & Role & Price (€) \\
\midrule
MCU & Control & 2.10 \\
LDO & Power & 0.35 \\
\bottomrule
\end{tabular}
\end{table}

\begin{table}
\begin{tabular}{lcr}
\hline
Left & Centre & Right \\
\hline
a & b & c \\
\hline
\end{tabular}
\caption{Column alignment: \texttt{l}, \texttt{c} and \texttt{r}}
\label{tab:align}
\end{table}

\subsection{Charts}

\begin{figure}
\begin{barchart}[ylabel={Settling time (ms)}, series={before, after}]
Step   & 42 & 18 \\
Ramp   & 35 & 21 \\
Impulse & 28 & 12 \\
\end{barchart}
\caption{A bar chart from a small table of numbers}
\label{fig:results}
\end{figure}

\subsection{References}

\texttt{\textbackslash{}ref} is the number alone, as in Table~\ref{tab:parts}. \texttt{\textbackslash{}Cref} and \texttt{\textbackslash{}autoref} add the word: \Cref{fig:design}, \autoref{tab:align}. \texttt{\textbackslash{}cref} gives it in lower case, \cref{fig:results}, and \texttt{\textbackslash{}eqref} an equation in brackets, \eqref{eq:pi}. Sections can be labelled too: see \Cref{sec:code}.

\clearpage

\section{Citations and bibliography}
\label{sec:cite}

Cite sources with \texttt{\textbackslash{}cite}, as in \cite{datasheet}; two at once give \cite{datasheet,lecture}, and a page goes in brackets: \cite[p.~12]{datasheet}. Numbers follow the order of the \texttt{\textbackslash{}bibitem}s, and the list is printed where \texttt{thebibliography} stands. A key that is neither a label nor a bibitem is printed as a visible unresolved reference, so look for one before publishing.

\begin{thebibliography}{9}
\bibitem{datasheet} Manufacturer, \emph{Part Datasheet}, rev.~3, 2025.
\bibitem{lecture} A.~Teacher, \emph{Control Systems}, lecture notes, 2024.
\end{thebibliography}

\section*{Acknowledgements}

\texttt{\textbackslash{}section*} gives an unnumbered heading that stays on the page it is written on, like this one.

% \appendix starts the appendices: their sections are lettered A, B, … and the
% page is called "Appendices". Delete this part if the post has none.
\clearpage
\appendix

\section{Wiring}
\label{app:wiring}

Material that would interrupt the argument, such as the full wiring table, goes here and is cited from the text as \Cref{app:wiring}.

\end{document}
